% ============================================================
% Appendix: Formal Proofs
% ============================================================

\section{Formal Proofs}
\label{sec:proofs}

We provide complete proofs for the four theorems stated in
\S\ref{sec:theorems}.  We also analyze the Cantor addressing scheme.

% ------------------------------------------------------------
\subsection{Cantor Pairing Properties}
\label{sec:cantor-analysis}

\begin{lemma}[Cantor Pairing Bijectivity]
\label{lem:cantor-bijection}
The standard Cantor pairing function
$\pi: \Nat \times \Nat \to \Nat$ defined by
$\pi(a, b) = \frac{(a+b)(a+b+1)}{2} + b$
is a bijection.
\end{lemma}

\begin{proof}
The function $\pi$ enumerates the pairs $(a,b)$ along successive
anti-diagonals $a+b = 0, 1, 2, \dots$.  On diagonal $n = a+b$,
the pairs are $(n,0), (n{-}1,1), \dots, (0,n)$, which are $n+1$
pairs.  The total number of pairs on diagonals $0, \dots, n{-}1$ is
$\sum_{i=0}^{n-1}(i+1) = \frac{n(n+1)}{2}$.  Thus
$\pi(a,b) = \frac{(a+b)(a+b+1)}{2} + b$ assigns a unique natural
number to each pair, and every natural number is reached.

The inverse is computed by finding the diagonal
$n = \lfloor\frac{\sqrt{8z+1}-1}{2}\rfloor$ for a given $z = \pi(a,b)$,
then $b = z - \frac{n(n+1)}{2}$ and $a = n - b$.
\end{proof}

\begin{lemma}[Iterated Cantor Pairing]
\label{lem:iterated-cantor}
The iterated (left-fold) Cantor pairing
$\cantor: \Nat^d \to \Nat$ defined by
$\cantor(p_1, \dots, p_d) = \pi(\pi(\cdots\pi(\pi(p_1, p_2), p_3)\cdots), p_d)$
is injective for each fixed depth $d$.
\end{lemma}

\begin{proof}
By induction on $d$.  Base case $d=2$: $\cantor(p_1, p_2) = \pi(p_1, p_2)$
is bijective by Lemma~\ref{lem:cantor-bijection}.  Inductive step:
assume $\cantor_d: \Nat^d \to \Nat$ is injective.  Then
$\cantor_{d+1}(p_1, \dots, p_{d+1}) = \pi(\cantor_d(p_1, \dots, p_d), p_{d+1})$.
Since $\pi$ is bijective and $\cantor_d$ is injective, the composition
is injective.  The iterated pairing is injective but not surjective
onto $\Nat$ for $d>2$ (not all naturals appear as valid iterated
encodings, since the first argument of $\pi$ is restricted to the
range of $\cantor_d$).
\end{proof}

\begin{remark}[Depth Truncation]
Our implementation truncates at $d_{\max} = 8$, bounding the Cantor
address to at most $\pi^{(7)}$ nested applications.  With branching
factor $k$, the maximum address value is $\cantor(k{-}1, k{-}1, \dots, k{-}1)$
for $d_{\max}$ components.  For $k=3$, $d_{\max}=8$, this produces
addresses up to $\sim 10^{12}$, well within 64-bit integer range.
The slot mapping $h(v) = \cantor(v) \bmod \slots$ may produce
collisions, which are resolved by open addressing with linear probing
in the cache.
\end{remark}


% ------------------------------------------------------------
\subsection{Proof of Theorem~\ref{thm:bounded} (Bounded Memory)}
\label{sec:proof-bounded}

\begin{proof}
The BIC system has two memory components:

\emph{(i) Cache array.}  The cache is a fixed-size array of $\slots$
slots.  Each slot stores at most one agent state of size
$\leq C_{\max}$.  Therefore the cache uses at most
$\slots \cdot C_{\max}$ bytes, regardless of $N$.

\emph{(ii) Agent registry.}  Each agent ever spawned gets a registry
entry containing: agent ID (8 bytes), parent ID (8 bytes), tree depth
(4 bytes), child indices (up to $k \times 4$ bytes), status (1 byte),
Cantor address (8 bytes).  For branching factor $k$:
$|r| \leq 29 + 4k$ bytes per agent.  With $k=3$, this is 41 bytes.
After $N$ spawn events, the registry uses $N \cdot |r|$ bytes, i.e.,
$\Oh(N)$.

\emph{(iii) No other state.}  The eviction manager, Hilbert compressor,
and API layer use $\Oh(\slots)$ auxiliary memory (priority queues,
index arrays).

Total: $\slots \cdot C_{\max} + N \cdot |r| + \Oh(\slots)
= \slots \cdot C_{\max} + \Oh(N)$.

Crucially, $\slots$ is a fixed constant chosen at initialization.
The $\Oh(N)$ registry term grows linearly, but each record is
$\sim$$|r| \approx 41$ bytes vs.\ $C_{\max} \approx 1{,}000$--$4{,}000$
bytes per state, so the registry is $\sim$$40\times$ smaller per agent
than the cache.
\end{proof}


% ------------------------------------------------------------
\subsection{Proof of Theorem~\ref{thm:frag} (Fragmentation)}
\label{sec:proof-frag}

\begin{proof}
The \textsc{Compact} operation is invoked explicitly (e.g., after a
batch of evictions) to restore contiguity; it is \emph{not} called on
every eviction, since the cache uses open addressing with linear probing
which tolerates sparse layouts.

The operation works as follows: given the cache array
$A[0..\slots{-}1]$ with occupied slots at arbitrary positions, we
scan left-to-right with two pointers: $w$ (write) starting at 0,
and $r$ (read) starting at 0.

\begin{enumerate}
  \item \textbf{Invariant:} At the start of each iteration,
    $A[0], A[1], \dots, A[w{-}1]$ are all occupied.
  \item While $r < \slots$:
    \begin{itemize}
      \item If $A[r]$ is occupied: set $A[w] \gets A[r]$ (move entry),
            update the slot-mapping index from $r \to w$,
            increment both $w$ and $r$.
      \item If $A[r]$ is empty: increment $r$ only.
    \end{itemize}
  \item Clear all slots $A[w..\slots{-}1]$.
\end{enumerate}

After compaction, all $|\BIC|$ occupied entries (where $|\BIC| = w$)
reside in contiguous positions $A[0..w{-}1]$.  The invariant is
maintained: $w$ only increments when an entry is placed, so
$A[0], A[1], \dots, A[w{-}1]$ are all occupied.

Positions $A[w..\slots{-}1]$ are all empty.  Therefore
\emph{external fragmentation is zero}: there are no empty slots
between occupied ones.

The operation runs in $\Oh(\slots)$ time (single pass over the array)
and requires no auxiliary memory beyond the two pointers.
\end{proof}


% ------------------------------------------------------------
\subsection{Proof of Theorem~\ref{thm:retrieve} (Retrievability)}
\label{sec:proof-retrieve}

\begin{proof}
We prove both parts.

\emph{Part (a): Cached agents.}  If agent $v$ has a cache entry at
slot $h(v) \bmod \slots$, the query resolves by hash-map lookup in
$\Oh(1)$ expected time (with open addressing and load factor
$< 1$, by standard hashing analysis).  The returned state is the
exact stored state.

\emph{Part (b): Evicted agents.}  If agent $v$ was evicted, its state
was summarized into its parent's entry via $\summarize$ at eviction
time.  To reconstruct $v$'s state, we walk the ancestor chain:

Let $v = v_d \to v_{d-1} \to \cdots \to v_1 \to v_0$ be the path from
$v$ to the root.  Define $v_{j^*}$ as the deepest ancestor that is
still cached (i.e., $j^* = \max\{j : v_j \in \mathrm{Cache}\}$).
Such $j^*$ exists because the root $v_0$ is never evicted (it has the
shallowest depth, so it is the last candidate for eviction; and at
cache size $\slots \geq 1$, at least one agent is always cached).

The reconstruction walks from $v$'s registry entry upward until
finding $v_{j^*}$, collecting summaries stored at $v_{j^*}$'s
cache entry.  This walk visits $d - j^*$ registry entries, each
costing $\Oh(1)$ (hash-map lookup), so the total cost is $\Oh(d)$.

\emph{Information loss bound.}  Each eviction-summarization step
preserves at least a $\preserv$ fraction of information.  An agent
at depth $d$ has been summarized at most $d - j^*$ times (once per
eviction along its ancestor chain).  In the worst case ($j^* = 0$,
ancestors evicted in order from leaf to root, then re-cached), the
cumulative preservation is $\preserv^{d}$.  Information loss is
bounded by $1 - \preserv^{d}$.

For $\preserv = 1 - \epsilon$ with small $\epsilon$:
$\preserv^d = (1-\epsilon)^d \geq 1 - d\epsilon$
(by Bernoulli's inequality).  With branching factor $k$, the maximum
depth is $d \leq \log_k N$, so reconstruction cost is
$\Oh(\log_k N)$ and information loss is at most
$1 - (1-\epsilon)^{\log_k N}$.
\end{proof}


% ------------------------------------------------------------
\subsection{Proof of Theorem~\ref{thm:liveness} (Indefinite Liveness)}
\label{sec:proof-liveness}

\begin{proof}
We must show that every \textsc{Spawn} operation completes
successfully, regardless of the number of agents previously spawned.

\emph{Case 1: Cache has a free slot.}
The new agent is placed in the free slot in $\Oh(1)$ time.

\emph{Case 2: Cache is full ($|\BIC| = \slots$).}
\textsc{EnsureCapacity} is invoked.  The eviction manager selects
at least one victim agent $v^*$ from the cache using the priority
scheme (idle $\to$ depth $\to$ LRU).  Since $|\BIC| = \slots \geq 1$,
at least one agent exists in the cache, so the selection always
succeeds.

The eviction step:
\begin{enumerate}
  \item Summarizes $v^*$'s state into its parent (if parent is cached).
  \item Removes $v^*$ from the cache, freeing one slot.
\end{enumerate}

After eviction, $|\BIC| = \slots - 1 < \slots$, so the new agent
can be placed.

\emph{Time complexity.}  Each operation's worst-case cost:
\begin{itemize}
  \item \textsc{Spawn}: $\Oh(\slots)$ (eviction candidate selection
    scans the cache).
  \item \textsc{Execute}: $\Oh(C_{\max})$ (state update).
  \item \textsc{Query}: $\Oh(d) \leq \Oh(\log_k N)$ by
    Theorem~\ref{thm:retrieve}.  Since $d \leq d_{\max} = 8$,
    this is $\Oh(1)$ in practice.
  \item \textsc{Compact}: $\Oh(\slots)$ (single array scan).
\end{itemize}

All operations complete in $\Oh(\slots)$ time, independent of $N$.
Since $\slots$ is finite, the system never blocks and can execute
indefinitely.

\emph{Liveness argument.}  No operation can deadlock: there are no
locks, no circular waits, and no unbounded loops.  The eviction
manager always makes progress (removes at least one agent), so the
system is live.
\end{proof}


% ------------------------------------------------------------
\subsection{Hilbert Curve Locality}
\label{sec:hilbert-analysis}

\begin{lemma}[Hilbert Locality Bound]
\label{lem:hilbert-locality}
For the $p$-order Hilbert curve in $n$ dimensions, two points at
$\ell_\infty$ distance $\leq r$ in $[0, 2^p)^n$ have Hilbert index
distance at most $(2r+1)^n \cdot 2^{n-1}$.
\end{lemma}

\begin{proof}[Proof sketch]
The Hilbert curve of order $p$ partitions $[0, 2^p)^n$ into $2^{np}$
cells, each a unit hypercube.  Two cells adjacent in
$\ell_\infty$ distance $\leq r$ lie within a hypercube of side
$2r+1$.  This hypercube contains at most $(2r+1)^n$ cells.  The
Hilbert curve visits each cell exactly once, and within any
axis-aligned sub-hypercube of side $s$, the curve's path has at most
$s^n \cdot 2^{n-1}$ local pieces (by the recursive
construction of the Hilbert curve).  Therefore the maximum index gap
between two cells in the sub-hypercube is bounded by
$(2r+1)^n \cdot 2^{n-1}$.

A full proof follows from the recursive structure of the Hilbert curve
and the Butz--Moore characterization of locality
\citep{butz1969hilbert, moon2001hilbert}.
\end{proof}

This locality property ensures that our Hilbert-based eviction
clustering co-evicts agents with similar state vectors, producing
more coherent summaries when the state space has spatial structure.


% ------------------------------------------------------------
\section{Coding-Task Case-Study Design Sketch}
\label{app:coding-case-study}

We sketch the next planned case study to demonstrate BIC's
deployment value on a workload structurally distinct from
research-decomposition.  No results are reported; this is a design
specification.

\paragraph{Task.}
A multi-agent coding pipeline that recursively converts a natural-language
specification into an executable Python module with passing unit
tests.  The orchestration layer is LangGraph~\citep{langgraph2024}
with Gemini-2.0-flash as the per-agent backend, sitting on top of
BIC for working-state management.

\paragraph{Sub-task tree.}
The root agent receives a specification (e.g., ``implement a
disjoint-set data structure with union-by-rank and path
compression'') and recursively decomposes:
\textbf{(1)~Parse spec} (root, depth 0): extract function signatures,
preconditions, and observable behaviours;
\textbf{(2)~Generate function stubs} (depth 1, $k=3$ children): one
agent per top-level function/method;
\textbf{(3)~Fill in implementations} (depth 2, $k=3$ children per
stub): one agent per implementation strategy
(naive / optimised / property-test scaffold);
\textbf{(4)~Run tests} (depth 3, $k=2$ children per implementation):
unit-test executor and property-based test executor;
\textbf{(5)~Debug} (depth 4, on-demand): spawned only when a test fails,
narrowing into the failing case.
Total tree size: $\sim$30--80 agents depending on debug depth.

\paragraph{Metrics.}
\textbf{(i)~test-pass rate}: fraction of generated unit tests that
pass on the final assembled module;
\textbf{(ii)~lines of correct code per dollar}: total accepted
implementation lines divided by aggregate Gemini-2.0-flash spend;
\textbf{(iii)~debug-depth distribution}: how often the tree extends
past depth 3, as a proxy for how aggressive cache pressure becomes
under realistic feedback loops;
\textbf{(iv)~task-coherence rate}: as in \S\ref{sec:case-study},
the fraction of sub-agent contributions retrievable by the root at
final assembly time.

\paragraph{Cache configurations to evaluate.}
We will sweep $\slots \in \{4, 8, 16, 32\}$ and the same backends
as \S\ref{sec:llm} (BIC, LRU, LRU+Summary, Unbounded).  We expect
BIC's structural advantage to be at least as pronounced as on
research-decomposition because debug spawns produce deeper trees
($d \geq 4$), exactly the regime where Theorem~\ref{thm:retrieve}'s
$\Oh(\log_k N)$ ancestor-chain reconstruction matters most.

\paragraph{Instrumentation.}
We will log per-eviction provenance (Cantor address, depth, summary
length) and per-agent test-pass outcomes.  Raw logs and the LangGraph
graph definition will be released alongside the camera-ready.


% ------------------------------------------------------------
\section{Failure Trace: Depth-3 Quality Dip Example}
\label{app:failure-trace}

To make the depth-decay failure mode of \S\ref{sec:failure-modes}
concrete, we extract one trace from the 50-agent LLM experiment
(\S\ref{sec:llm}) where BIC's semantic-reconstruction quality dipped
below the run mean.  All numbers come from the raw log file
\texttt{experiment\_6\_llm\_50agents\_raw.json} in the public
repository.

\paragraph{Configuration.}
Backend BIC, $\slots = 32$, seed 2024, branching factor $k=3$,
depth $d=3$.  Question: ``Survey the state of memory management in
multi-agent LLM systems''.  Total agents in this run: 40.

\paragraph{Aggregate metrics for this seed.}
Semantic reconstruction quality $0.228$ (vs.\ run mean $0.253$ at
$\slots{=}32$ over 5 seeds, range $[0.228, 0.276]$).  Query-success
rate $1.000$.  Cache-hit rate $\sim$0.32.  Total evictions during
this run: $\sim$56.  Information-preservation ratio $0.69$, in the
expected range for the default concat-merge summariser.

\paragraph{Where the quality came from.}
Inspecting per-agent reconstruction scores, the seed's overall mean
is pulled down by two depth-3 leaf agents (the deepest agents in this
$d=3$ tree) whose reconstructions retained the parent's framing
tokens (e.g.\ ``memory management'', ``LLM agents'') but lost the
concrete domain-specific vocabulary in their original answers
(e.g.\ specific paper titles, system names, numeric thresholds).
TF-IDF cosine similarity penalises this lexical loss heavily because
high-IDF concrete nouns dominate the score.  Query-success was
unaffected --- every leaf was retrievable through its ancestor chain
--- so the failure is one of \emph{fidelity}, not \emph{availability}.

\paragraph{Mechanism.}
The two leaves shared a depth-2 grandparent that itself was evicted
mid-run.  The grandparent's summary was folded into its own parent
(depth 1) before either leaf was queried by the reconstruction
walker.  The leaves' content therefore traversed two
summarisation steps before reaching a cached ancestor; with
$\preserv \approx 0.74$ per step, the cumulative bound is
$0.74^2 \approx 0.55$, consistent with the depth-decay model.

\paragraph{Mitigation suggested by this trace.}
Either (i)~increase $\slots$ enough to keep all depth-2 grandparents
cached for the duration of the run --- which on this $N=40$ tree
requires $\slots \geq 13$ in the worst case, comfortably below the
$\slots{=}32$ used here, but only if the eviction priority is biased
against intermediate-depth agents; or (ii)~switch to the LLM-backed
$\summarize$, which we hypothesise would preserve more of the
high-IDF concrete vocabulary and therefore close the depth-3 gap.
We have not yet run the LLM-backed configuration; we mark this as
the most-actionable item for the camera-ready evaluation.

% TODO: Karim, add additional concrete traces if available --- a deeper
% tree ($d \geq 5$) example would be the strongest illustration of the
% depth-decay mode.
